\documentclass[11pt,letterpaper]{article}

\usepackage[margin=1in]{geometry}
\usepackage{times}
\usepackage{amsmath,amssymb}
\usepackage{booktabs}
\usepackage[colorlinks=true,linkcolor=black,citecolor=blue,urlcolor=blue]{hyperref}
\usepackage{caption}
\captionsetup{font=small,labelfont=bf}

\title{\bfseries Frequency-Aware DETR Training at Extreme Class Imbalance}
\author{Technical report}
\date{}

\begin{document}
\maketitle

\begin{abstract}
\noindent
Logit adjustment is a standard remedy for long-tailed recognition, and applying
it consistently across a DETR-family model's classification loss, Hungarian
matching cost and denoising task is a natural extension. We report that at
34{,}320:1 imbalance it fails catastrophically: the unified treatment is
6.05\,pp mAP below a stock baseline with tail AP at exactly zero, and loses
4.41\,pp to conventional class-balanced reweighting. We identify the cause as
magnitude rather than principle---at $\tau=1.0$ the adjustment applies a
$+11.47$ logit shift to the rarest class---and show recovery is monotonic in
$\tau$ while never clearing the baseline. We further report a reproducible
stability pattern: every arm applying the adjustment \emph{inconsistently}
diverged, while every consistent arm trained stably. Finally, decoupled
classifier retraining is nominally $+0.53$\,pp mAP on test, of which 98\,\%
comes from a single class with \emph{two} test instances.
\end{abstract}

\section{Setup}

Stock DINO-DETR~\cite{zhang2023dino}, ResNet-50 4-scale, initialised from the
official COCO checkpoint with a reinitialised 31-class embedding.
$\mathrm{num\_classes}=32$ because DINO indexes raw \texttt{category\_id}
without remapping.

All arms: 12 epochs (official $1\times$), \texttt{lr\_drop}~11, batch~2,
seed~42, identical initialisation and evaluation protocol. \textbf{No arm
receives a hyperparameter search.} That keeps budgets matched; Section~4 shows
it is also, in this instance, the defect.

Evaluated on a verified-disjoint split: 9752 / 2090 / 2090 images, 0 exact
duplicates and 0 NCC-confirmed near-duplicates across 792{,}811 candidate pairs.

\section{The Claim Under Test}

In DETR-family models the matcher decides which query is supervised for which
ground truth. If the classification loss is frequency-adjusted but the matching
cost is not, matching keeps assigning tail ground truth to queries whose
unadjusted scores look good while the loss optimises a different scale. The
hypothesis is that applying the adjustment \emph{consistently}---across loss,
matching cost and denoising---beats applying it in one place, and beats ordinary
loss reweighting.

The matrix was frozen before any arm ran. The control C1 applies per-class loss
weights only~\cite{cui2019classbalanced}---no logit shift, nothing in the
matcher or denoising---and was verified to reproduce stock sigmoid focal loss to
\textbf{0.0 absolute difference} at unit weights, so the per-class weight is
provably the only change.

\section{Result}

\begin{table}[h]
\centering\small
\begin{tabular}{llrrrrr}
\toprule
arm & configuration & mAP & AP50 & AP75 & mid & tail \\
\midrule
D1 & baseline & \textbf{0.1625} & 0.3080 & 0.1551 & \textbf{0.2118} & 0.0368 \\
D2 & loss only & \multicolumn{5}{c}{\emph{diverged}} \\
D3 & matching only & \multicolumn{5}{c}{\emph{diverged}} \\
D4 & denoising only & 0.1633 & \textbf{0.3165} & \textbf{0.1572} & 0.2029 & \textbf{0.0455} \\
D5 & \textbf{unified} & 0.1020 & 0.1998 & 0.0939 & 0.0959 & \textbf{0.0000} \\
D6 & unified $+$ oversampling & 0.1014 & 0.1988 & 0.0946 & 0.0943 & 0.0000 \\
D7 & unified $+$ contrast & 0.1012 & 0.1993 & 0.0929 & 0.0972 & 0.0000 \\
C1 & plain reweighting & 0.1461 & 0.2910 & 0.1375 & 0.1847 & 0.0270 \\
\bottomrule
\end{tabular}
\caption{Validation, matched budget. The decisive comparison $D5-C1$ is
mAP $-4.41$, AP50 $-9.12$, AP75 $-4.36$, mid $-8.87$, tail $-2.70$\,pp.}
\end{table}

The unified treatment loses decisively to plain class-balanced reweighting,
which in turn loses 1.64\,pp to changing nothing. \textbf{Tail AP in every
unified arm is exactly 0.0000}---the metric the method targets is the one it
destroys. Layering oversampling~\cite{gupta2019lvis} or contrast enhancement
changes nothing, because the base configuration is already broken.

Only D4 clears the baseline, by 0.08\,pp mAP. That is inside noise and carries
no claim. On held-out test, D5 confirms: mAP 0.1001 against 0.1570
($-5.69$\,pp), tail 0.0000 against 0.0598.

\section{Diagnosis: Magnitude, Not Principle}

Logit adjustment~\cite{menon2021logit} subtracts $\tau\log p(c)$ from each class
logit. On this training set the rarest present class ($n=1$) has log-prior
$-11.47$ and the most common ($n=34{,}318$) has $-1.03$: a spread of
\textbf{10.44} in logit space. A $+11.47$ shift saturates the sigmoid for every
query on that class.

Menon~et~al. validated $\tau=1.0$ on CIFAR-LT and ImageNet-LT, whose imbalance
is of order 100:1 to 1000:1. This corpus is \textbf{34{,}320:1}.

\begin{table}[h]
\centering\small
\begin{tabular}{rrrr}
\toprule
$\tau$ & mAP & $\Delta$ vs baseline & tail \\
\midrule
1.00 & 0.1020 & $-6.05$ & 0.0000 \\
0.50 & 0.1340 & $-2.85$ & 0.0135 \\
0.25 & 0.1601 & $-0.24$ & 0.0346 \\
\bottomrule
\end{tabular}
\caption{Sweep on the unified configuration, changing nothing else. Recovery is
monotonic in $\tau$, confirming magnitude as the mechanism, and incomplete:
even at $\tau=0.25$ the method does not reach the baseline.}
\end{table}

This is a transferability result about the published default. $\tau=1.0$ is used
across the long-tail literature without restatement of the imbalance regime it
was tuned for. At clinical imbalance it is not merely suboptimal; it is
destructive.

\section{Divergence Is Patterned}

Four arms failed with \texttt{inf} or \texttt{NaN} cost matrices at
\texttt{linear\_sum\_assignment}. Which arms failed is not random: D2 (loss
only), D3 (matching only), L1 (matching $+$ denoising) and L2 (loss $+$
denoising) all diverged, while D4 (denoising only) and D5--D7 (all three) were
stable.

Checkpoints on disk are numerically clean---max $|w|$ $1.42\mathrm{e}{+}01$,
identical to arms that completed---so divergence develops during the run.
Gradient clipping was active at 0.1 throughout. When the loss carries a $+11.47$
shift and the matcher scores on unadjusted probabilities, the two components
optimise against different scales and the assignment degenerates.
\textbf{Consistency prevents divergence; it does not produce accuracy.}

We report this as an observation from crashes, not a designed stability
experiment. The pattern is reproducible---D2 diverged on every attempt.

\section{Decoupled Classifier Retraining}

The long-tail literature indicates rebalancing belongs at the classifier
stage~\cite{kang2020decoupling,wang2019calibration}, not inside representation
learning---the opposite of every arm above, and structurally unable to exhibit
the failure of Section~5. We freeze all parameters except the classification
embedding, start from the trained baseline, and retrain for 6 epochs under
repeat-factor sampling.

On test this gives mAP 0.1623 against 0.1570 ($+0.53$\,pp), AP50 $+0.88$\,pp,
tail $+1.25$\,pp---positive on every metric, with the largest gain exactly where
a long-tail method should deliver.

\textbf{It does not survive decomposition.} The largest single contributor is
one class with \textbf{two test instances} moving $0.0000\to0.1515$,
contributing $+0.52$\,pp---\textbf{98\,\% of the total}. Excluding it, the gain
is $+0.01$\,pp. On the 18 classes with ${\ge}10$ test instances, 6 improve and
12 worsen.


\section{Noise Floor (Measured)}

The reference configuration was trained at three seeds, everything else
identical. Determinism is exact on this pipeline, so seed choice is provably the
only source of variation.

\begin{table}[h]
\centering\small
\begin{tabular}{lrrrrr}
\toprule
seed & mAP & AP50 & AP75 & head & tail \\
\midrule
42 & 0.1055 & 0.2549 & 0.0661 & 0.2815 & 0.0101 \\
1337 & 0.1056 & 0.2573 & 0.0679 & 0.2837 & 0.0117 \\
2024 & 0.1037 & 0.2568 & 0.0629 & 0.2823 & 0.0105 \\
\midrule
\textbf{sd} & \textbf{0.0010} & 0.0013 & \textbf{0.0025} & 0.0011 & 0.0008 \\
\bottomrule
\end{tabular}
\caption{Noise floor at 2\,sd: mAP $\pm0.21$, AP50 $\pm0.26$,
AP75 $\pm0.50$, head $\pm0.22$, tail $\pm0.17$\,pp.}
\end{table}

Three consequences, and they are not the same verdict. The seven-arm ablation
spans a 0.36\,pp range, \emph{inside} the band---selecting a winner on mAP was
selecting noise. The best segmentation arm on test is $+0.21$\,pp, sitting
\emph{exactly at} the threshold. But the isolated boundary effect on AP75 is
$+0.70$\,pp against a $\pm0.50$\,pp floor, so it \textbf{exceeds} noise on
validation at 50 epochs, and then reverses on test at 100 epochs. That is a
\textbf{transfer failure, not a noise result}: the objective does what its
mechanism predicts under the conditions it was measured in, and that does not
survive a change of split and schedule.

\section{Conclusion and Limitations}

No configuration tested improves on stock DINO-DETR beyond noise. Consistent
frequency-awareness loses 4.41\,pp to plain reweighting, which loses to changing
nothing; the pre-registered hypothesis is refuted by its own frozen matrix. The
published $\tau=1.0$ does not transfer to clinical imbalance, with failure
monotonic in $\tau$. Inconsistent application diverges reproducibly. Decoupled
classifier retraining is the right mechanism and still does not clear noise. And
group means over low-support classes fabricate results---twice, in two
independent projects.

Single seed per arm, with the noise floor measured above. Single corpus.
No external benchmark, and therefore no long-tail claim in the title, abstract
or conclusions. Two cells are recorded as diverged rather than scored.

{\small
\bibliographystyle{plain}
\bibliography{references}
}

\end{document}
